\section{Proposed Method}
\label{sec:method}

This section formalises the steppable-surface segmentation problem. It then
describes a parametric generator that produces a fully labelled staircase point
cloud with realistic sensor artefacts. Finally, it details four classical
geometric methods that we implement and compare under an identical 3-class
formulation.

\subsection{Problem Formulation and Class Scheme}
We treat staircase perception as a point-wise classification problem. Given an
unorganised cloud $\mathcal{P}=\{\mathbf{x}_i\}_{i=1}^{M}$ with
$\mathbf{x}_i\in\R^3$, the goal is to assign every point a label
$\ell_i\in\{0,1,2\}$, where label $0$ denotes a \emph{tread} (a steppable,
nominally horizontal surface), label $1$ denotes a \emph{riser} (the vertical
face connecting consecutive treads), and label $2$ denotes \emph{other}
(structural edges, sensor noise, and outliers). The tread class is the primary
class of interest, since a robot or assistive device plans footholds on tread
surfaces. We therefore evaluate it in a one-vs-rest setting while still reporting
overall 3-class accuracy.

The decisive geometric cue separating the first two classes is surface
orientation relative to the vertical axis. Let $\nvec$ be the unit surface
normal at a point and $\zhat$ the world vertical. A tread satisfies
$|\nvec^\top\zhat|\to 1$ (its normal is nearly parallel to $\zhat$), whereas a
riser satisfies $|\nvec^\top\zhat|\to 0$ (its normal is nearly orthogonal to
$\zhat$). Every method below exploits this orientation contrast, either
explicitly through a normal-based gate or implicitly through the height
structure of the cloud. Treating orientation as the organising principle yields
the unified pipeline summarised in Algorithm~\ref{alg:pipeline}. Each of the four
methods plugs into this pipeline at the clustering or plane-fitting stage.

\begin{algorithm}[t]
\caption{Orientation-aware segmentation (generic)}\label{alg:pipeline}
\begin{algorithmic}[1]
\STATE estimate per-point normals $\nvec_p$ via PCA over $k$ neighbours \eqref{eq:cov}
\STATE split points into horizontal / vertical orientation groups
\STATE cluster or fit planes within each group
\STATE confirm each region by its plane tilt $|\nvec^\top\zhat|$
\STATE label tread / riser / other
\end{algorithmic}
\end{algorithm}

\subsection{Parametric Data Generation}
To obtain a controlled and fully labelled benchmark, we synthesise the staircase
analytically from physical parameters: width $w=2.5$\,m, tread depth
$d=0.30$\,m, step height $h=0.18$\,m, and $N=4$ steps. For step index
$i\in\{0,\dots,N-1\}$, tread points are sampled uniformly over the horizontal
extent of step $i$ and placed at height $ih$ with additive Gaussian height
noise following \eqref{eq:tread}:

\begin{equation}\label{eq:tread}
\begin{cases}
x \sim \mathcal{U}(0,\,w) \\
y \sim \mathcal{U}(i d,\,(i{+}1) d) \\
z = i h + \epsilon_z,\quad \epsilon_z \sim \mathcal{N}(0,\sigma_z^2)
\end{cases}
\end{equation}
\begin{equation}\label{eq:riser}
\begin{cases}
x \sim \mathcal{U}(0,\,w) \\
y = i d + \epsilon_y,\quad \epsilon_y \sim \mathcal{N}(0,\sigma_y^2) \\
z \sim \mathcal{U}((i{-}1) h,\, i h)
\end{cases}
\end{equation}

Riser points come from \eqref{eq:riser} for $i\in\{1,\dots,N-1\}$, the vertical
face rising to tread $i$. They share the same lateral spread but lie on a
vertical plane at depth $id$ with small along-depth jitter $\sigma_y=0.02$\,m,
and their height sweeps uniformly across one step rise. The nominal height noise
is $\sigma_z=0.02$\,m, which we later vary to probe sensitivity. On top of this
structured cloud we inject realism in two ways. About $4\%$ of the points are
replaced by uniformly distributed outliers spanning the staircase bounding
volume, mimicking spurious returns and edge artefacts; these carry the
\emph{other} label. Every retained point then receives an isotropic LiDAR-like
jitter drawn from $\mathcal{N}(0,0.005^2)$ on each coordinate, approximating the
few-millimetre range noise typical of indoor depth sensors. We treat this value
as a nominal modelling choice, not a calibrated sensor model, and probe its
principal axis ($\sigma_z$) directly in Section~\ref{sec:evaluation}. The
resulting cloud contains $18{,}333$ points, partitioned as $12{,}800$ tread,
$4{,}800$ riser, and $733$ other. The generator emits per-point labels by
construction, so it provides exact ground truth for both training-free
evaluation and the noise sweep. This removes the manual-annotation bottleneck
that hampers learned segmenters.

\subsection{RANSAC Plane Fitting}
The first method fits planes by Random Sample Consensus. A candidate plane with
unit normal $\nvec$ and offset $b$ scores points by their orthogonal distance.
The inlier set $\mathcal{I}$ then collects every point closer than a tolerance
$\tau$:

\begin{equation}\label{eq:plane}
d(\mathbf{x}) = |\,\nvec^\top \mathbf{x} + b\,|,\qquad
\mathcal{I} = \{\mathbf{x} : d(\mathbf{x}) < \tau\}
\end{equation}

A naive single fit fails on stairs. The dominant plane that maximises inliers is
the oblique surface spanning the entire staircase, whose normal has
$|\nvec^\top\zhat| = d/\sqrt{d^2+h^2}\approx 0.86$. RANSAC therefore returns a
spurious ``ramp'' that absorbs both treads and risers into one steppable region.
We avoid this with a two-phase, orientation-gated scheme. The first phase
extracts horizontal planes (treads) by accepting only fits whose normal
satisfies $|\nvec^\top\zhat|\ge 0.93$. The second re-runs RANSAC on the residual
points to extract vertical planes (risers) with $|\nvec^\top\zhat|\le 0.25$. We
set $\tau=0.035$\,m and impose minimum inlier counts of $2{,}000$ points for
tread planes and $300$ for riser planes. This suppresses thin riser slices from
contaminating tread fits, at the cost of occasionally rejecting sparse vertical
faces. The orientation-aware variant follows the spirit of structure-guided
RANSAC refinements~\cite{li2017improved,jalal2021rgb}.

\subsection{PCA Normal Analysis}
The second method classifies each point purely from its local surface
orientation. For point $p$ we gather its $k$ nearest neighbours
$\mathcal{N}_k(p)$ and form the local covariance matrix $\mathbf{C}_p$. The
eigenvector tied to its smallest eigenvalue serves as the surface normal:

\begin{equation}\label{eq:cov}
\mathbf{C}_p = \frac{1}{k}\sum_{j\in\mathcal{N}_k(p)}
(\mathbf{x}_j-\bar{\mathbf{x}})(\mathbf{x}_j-\bar{\mathbf{x}})^\top,\quad
\nvec_p = \mathbf{v}_{\min}(\mathbf{C}_p)
\end{equation}

The smallest-eigenvalue direction is the axis of least point spread, which for a
locally planar patch coincides with its normal~\cite{huang2009consolidation,castillo2013point}.
With $k=50$ the estimate stays stable against the injected jitter while remaining
local enough to respect tread and riser boundaries. We classify by the
verticality of the estimated normal: a point is labelled tread when
$|\nvec_p^\top\zhat|\ge 0.80$ and riser when $|\nvec_p^\top\zhat|\le 0.50$.
Intermediate values, which arise mostly at edges and from outliers, fall to the
\emph{other} class. This decision is per-point and parameter-light, with no
plane-fitting iterations. The PCA thresholds ($0.80/0.50$) are deliberately
looser than the RANSAC orientation gate ($0.93/0.25$). RANSAC tests a global
plane whose aggregate normal is precise, whereas per-point PCA normals are
noisier and need a wider acceptance band.

\subsection{DBSCAN with Slope Confirmation}
The third method couples orientation reasoning with spatial clustering. We first
pre-split the cloud into a horizontal group and a vertical group using the same
normal-verticality criterion as above. Within each group we run density-based
spatial clustering (DBSCAN) with neighbourhood radius $\varepsilon=0.08$\,m, so
that each physically separate surface forms its own connected cluster, while
sparse outliers, lacking dense support, fall to the noise label. Pre-splitting by
orientation matters here. Without it, a tread and the riser directly beneath it,
which are spatially adjacent, would merge into a single density-connected
cluster. Each resulting cluster is confirmed by its aggregate slope, that is the
verticality $|\nvec^\top\zhat|$ of its fitted plane, and labelled tread, riser, or
other accordingly. This pairing of clustering with region-level geometric
confirmation mirrors established point-cloud clustering and region-growing
practice~\cite{zhou2024identification,luo2021supervoxel}.

\subsection{Height-Histogram Analysis}
The fourth method is the most specialised to the staircase structure and ignores
normals entirely. Every tread of an ideal staircase lies at a discrete height
$ih$, so the marginal distribution of the $z$ coordinate concentrates into sharp
peaks at those levels, separated by low-density bands that correspond to the
riser sweeps. We build a one-dimensional histogram of $z$ over the whole cloud
and detect its dominant peaks. Points falling within a window of half-width
$0.05$\,m around each peak are labelled tread, points in the intervening height
bands are labelled riser, and the remainder become other. On the benchmark cloud
this recovers four tread levels at $z\approx\{0.00, 0.18, 0.36, 0.54\}$\,m,
matching the generative parameters. The method is extremely fast and needs no
neighbourhood computation, but it presumes a regular, axis-aligned staircase. As
the evaluation shows, its height bands tend to over-include riser points near
each tread level.

\subsection{Implementation and Parameters}
All methods were implemented in Python with NumPy and scikit-learn. RANSAC and
the normal-based gates are custom, whereas DBSCAN and the nearest-neighbour
search reuse standard library routines. Table~\ref{tab:params} consolidates every
parameter for reproducibility. The four methods share an asymptotic cost
dominated by the $k$-nearest-neighbour normal estimation ($O(M\log M)$ with a
$k$-d tree for $M$ points). On top of this, RANSAC adds a term linear in the
iteration budget, DBSCAN a near-linear neighbourhood query, and the
height-histogram stays at $O(M)$, which makes it by far the cheapest. All values
were held fixed across the experiments reported in
Section~\ref{sec:evaluation}.

\begin{table}[t]
\caption{Hyper-parameters held fixed across all experiments.}
\label{tab:params}
\centering
\footnotesize
\begin{tabular}{lll}
\toprule
Symbol & Value & Description \\
\midrule
$w,d,h$ & $2.5,\,0.30,\,0.18$\,m & width, tread depth, step height \\
$N$ & $4$ & number of steps \\
$\sigma_z,\sigma_y$ & $0.02,\,0.02$\,m & tread / riser Gaussian noise \\
jitter & $0.005$\,m & isotropic LiDAR-like jitter \\
outliers & $4\%$ & uniform outlier fraction \\
$\tau$ & $0.035$\,m & RANSAC inlier threshold \\
min.\ inliers & $2{,}000\,/\,300$ & RANSAC tread / riser plane \\
$k$ & $50$ & PCA normal neighbours \\
gate (RANSAC) & $0.93\,/\,0.25$ & horiz.\ / vert.\ $|\nvec^\top\zhat|$ \\
gate (PCA) & $0.80\,/\,0.50$ & tread / riser $|\nvec^\top\zhat|$ \\
$\varepsilon$ & $0.08$\,m & DBSCAN radius \\
window & $0.05$\,m & histogram peak half-width \\
\bottomrule
\end{tabular}
\end{table}
